\documentclass[11pt]{article}
\usepackage{fontspec}
\usepackage[french,english]{babel}
\usepackage{amsmath}
\usepackage{amssymb}
\usepackage{amsthm}
\usepackage{mathrsfs}
\usepackage{tikz-cd}
\usepackage{enumitem}
\usepackage{lmodern}

\newcommand{\printedpage}[1]{\marginpar{\footnotesize p.~#1}}
\newcommand{\msp}[1]{\marginpar{\footnotesize [#1]}}

\newcommand{\cE}{\mathcal{E}}
\newcommand{\cC}{\mathcal{C}}
\newcommand{\Ob}{\operatorname{Ob}}
\newcommand{\Fl}{\operatorname{Fl}}
\newcommand{\Hom}{\operatorname{Hom}}
\newcommand{\Ind}{\operatorname{Ind}}
\newcommand{\card}{\operatorname{card}}
\newcommand{\red}{\operatorname{\text{red}}}
\newcommand{\id}{\mathrm{id}}
\newcommand{\Ens}{\mathrm{Ens}}
\newcommand{\circled}[1]{\tikz[baseline=(c.base)]\node[draw,circle,inner sep=1pt](c){#1};}

\begin{document}

\printedpage{164}
% head: A. Grothendieck : Dérivateurs, Ch. XVIII, Édition des Universités de Montpellier II et Paris VII

\msp{page 430}to prove that for every $\lambda : s \longrightarrow t$, the arrow
\[
\varinjlim_{i} \Hom_{\lambda}(x, y_i) \ \longrightarrow \ \Hom_{\lambda}(x, y)
\]
is bijective. For fixed $x$ and $\lambda$, and varying direct systems in $\cE_t$, large relative to $\pi$, this means that the functor
\[
H_{\lambda}(x) \ = \ (y \longmapsto \Hom_{\lambda}(x, y)) \ : \ \cE_t \ \longrightarrow \ \Ens
\]
is $\pi$-accessible, and this for any
\[
x \ \in \ \Ob \cE_{s\,\pi} \ .
\]
Now this is the case if $s \in \Ob \cE_{s\,\pi_0}$, by virtue of hypothesis $3^{\circ}$) (then $H_{\lambda}(x_0)$ is $\pi_1$-accessible, % [sic] $s \in \Ob \cE_{s\,\pi_0}$ -> $x_0 \in \Ob \cE_{s\,\pi_0}$, the condition concerns an object of $\cE_{s\,\pi_0}$
and a fortiori $\pi$-accessible). We have to pass from there to the case where $x \in \Ob \cE_{s\,\pi}$. Now since $\pi_0$ is useful for $\cE_s$ by $1^{\circ}$, we have
\[
x \ = \ \varinjlim_{j \in J} x_j \ , \qquad x_j \in \Ob \cE_{s\,\pi_0},\ J \text{ large relative to } \pi_0 \ ,
\]
and since by $2^{\circ}$, $\pi^{\pi_0} = \pi \geqslant \pi_1 \geqslant \sup(\pi_0, \card \red \cE_{s\,\pi_0})$, we conclude from this that we can take $J$ so that
\[
\card J \ \leqslant \ \pi \ .
\]
On the other hand, by virtue of $\pi \geqslant \pi_1 \ \geqslant \ \pi_{\lambda}$, the functor
\[
\left\{
\begin{array}{lcl}
x_0 & \longmapsto & H_{\lambda}(x_0) \\[1ex]
\cE_s & \longrightarrow & (\underline{\Hom}(\cE_s, \Ens))^{\circ}
\end{array}
\right.
\]
is $\pi_0$-accessible, so we have
\[
H_{\lambda}(x_0) \ \simeq \ \varprojlim_{j \in J} H_{\lambda}(x_j) \ .
\]

\msp{page 431}Since $x_j \ \in \ \Ob \cE_{s\,\pi_0}$, $H(x_j)$ is a $\pi_1$-accessible functor, and since $\card J \leqslant \pi$, the $\varprojlim\limits_{j \in J} H_{\lambda}(x_j)$ is also $\pi$-accessible (cf. SGA 4 I 9.8).

Thus we have proved $\cE_{s\,\pi} \subset \cE(\pi)$ for every $s \in \Ob \cC$, i.e. the inclusion
\[
\cE(\pi) \ \subset \ \cE_{\pi} \ ,
\]
whence the fully faithful functor
\[
\Ind_{\pi}(\cE(\pi)) \longrightarrow \cE \ , \tag*{$(*)$}
\]
since $\cE$ satisfies $\mathrm{L}_{\pi}$. It remains to prove that it is essentially surjective. But if $x \in \Ob \cE$, $s = f(x)$, we have $x \in \Ob \cE_s$, and since $\pi_0$ is useful for $\cE_s$ by $1^{\circ}$), and $\pi^{\pi_0} = \pi > \pi_0$,
$\pi$ is also useful for $\pi$, so we have \textit{in} $\cE_s$ % [sic] utile pour $\pi$ -> utile pour $\cE_s$, usefulness is a property relative to the category $\cE_s$
\[
x \ = \ \varinjlim_{I} x_i \ , \qquad \text{the } x_i \in \cE_{s\,\pi} \subset \cE(\pi),\ I \text{ large relative to } \pi \ ,
\]
\printedpage{165}
% head: A. Grothendieck : Dérivateurs, Ch. XVIII, Édition des Universités de Montpellier II et Paris VII
and by a) (since $\pi \geqslant \pi_1 \geqslant \sup_{\lambda \in \Fl \cC} \pi_{\lambda}$) $x$ is also a limit of the $x_i$ in $\cE$, so it is in the essential image of $(*)$.

d) It remains to prove that we have $\cE_{\pi} \ \subset \ \cE(\pi)$. But we know (by 4.8.~\dots), since $\cE \ \simeq \ \Ind_{\pi}(\cE(\pi))$, that $\cE_{\pi} \ = \ \cE(\pi)^{\natural}$, i.e. that for $x \ \in \ \Ob \cE_{\pi}$, $\exists \ y \ \in \ \cE(\pi)$ and a projector $p$ in $x$, such that $x$ is isomorphic to the image of $p$ in $\cE$. Now $f(p)$ is a projector in $f(y) = s$, so by the hypothesis on $\cC$, $f(p) = \id_s$, i.e. $p$ is a projector in $\cE_s$. Since therefore $x$ is in $\cE_s$, and

\msp{page 432}even in $\cE_{s\,\pi}$, since $\cE_{s\,\pi}$ is stable under direct factors in $\cE_s$. This proves that $x \in \Ob \cE(\pi)$, q.e.d.

We can present 4.15.1 in a more congenial form.

\textbf{Corollary 4.15.4.} \textit{Let $f : \cE \longrightarrow \cC$ be a functor of $\mathfrak{U}$-categories. The following conditions are equivalent:}
\begin{enumerate}
\item[\textit{a)}] \textit{For every functor $\cC' \longrightarrow \cC$, with $\cC'$ small and without projectors, $\cE' \overset{\text{def}}{=} \cE \times_{\cC} \cC'$ is accessible, and $f' : \cE' \longrightarrow \cC'$ is accessible.}
\item[\textit{b)}] \textit{As in a), but with $\cC' = \mathbf{\Delta}_1$, in other words, this is condition $H$ p.~409: for every $\lambda \in \Fl \cC$, $\cE_{\lambda} = \cE \times_{\cC} (\mathbf{\Delta}_1, i_{\lambda})$ is accessible, and $\cE_{\lambda} \longrightarrow \mathbf{\Delta}_1$ is accessible.}
\item[\textit{c)}] \textit{The condition $H' = H\ \mathit{1} + H\ \mathit{2} + H\ \mathit{3}$, made explicit in the statement of 4.15.1.}
\end{enumerate}
\textsc{Proof}. It is done in a circle,
\[
\mathrm{a} \ \Longrightarrow \ \mathrm{b} \ \Longrightarrow \ \mathrm{c} \ \Longrightarrow \ \mathrm{a} \ ,
\]
the implication a $\Longrightarrow$ b being tautological, b $\Longrightarrow$ c being almost so (cf. 4.14.24, a $\Longrightarrow$ b $\Longrightarrow$ c). In fact, conditions b and c are conditions on the categories $\cE_s \xrightarrow{f_{\lambda}} \mathbf{\Delta}_1$ for $\lambda \in \Fl \cC$, hence stable under base change. To prove

\msp{page 433}c $\Longrightarrow$ a, we can therefore assume that $\cC' = \cC$, $\cC' \longrightarrow \cC$ the identity functor, and prove that if c) is satisfied and if $\cC$ is small and without projectors, then $\cE$ is accessible and $f$ accessible. Now this is what 4.15.1 asserts.

\textbf{Remark 4.15.5.} If $f : \cE \longrightarrow \cC$ is a fully faithful functor, it is immediate that condition H is satisfied (for $\cE_{\lambda} \longrightarrow \mathbf{\Delta}_1$ will be fully faithful, so $\cE_{\lambda}$ is equivalent
to $e$ or to $\mathbf{\Delta}_1$, hence accessible, and every functor from $\cE_{\lambda}$ to any category is % [sic] où -> ou, disjunction "e ou Δ1"
accessible \dots). But if we assume $\cC$ small \textit{without} assuming $\cC$ without projectors, it does not follow that $\cE$ is accessible, or what amounts to the same thing (being essentially small), that $\cE$ is Karoubian. (The most trivial example is the one where $\cE = \cC$, $f = \id_{\cC}$, $\cE$ is Karoubian only if $\cC$ is!) But even assuming $\cC$ a small Karoubian category, it does not necessarily follow that $\cE$ is. Indeed, \textit{every} small category $\cE$ (Karoubian or not) embeds as a full subcategory in its Karoubian envelope $\cC$ -- so it suffices to take $\cE$ non-Karoubian. The example

\msp{page 434}that is most economical is
\[
\Ob \cE \ = \ \{0\} \ , \qquad \Fl(\cE) \ = \ \{\id_0, p\} \ , \qquad \text{with } p^2 = p \neq \id_0 \ ,
\]
\printedpage{166}
% head: A. Grothendieck : Dérivateurs, Ch. XVIII, Édition des Universités de Montpellier II et Paris VII
the Karoubi envelope of $\cE$ is given by
\[
P \ = \ \left\{
\begin{tikzcd}
1 \arrow[r, shift right, "\sigma"'] & 0 \arrow[l, shift right, "\rho"'] \arrow[loop right, "p"]
\end{tikzcd}
\qquad
\begin{aligned}
\rho\sigma \ &= \ \id_1 \\
\sigma\rho \ &= \ p \ \neq \ \id_0
\end{aligned}
\quad \Longrightarrow \ p^2 = p
\right\} .
\]
This raises the question of finding a statement similar to 4.15.1, giving simple serviceable conditions, stronger than $\mathrm{H}' = \mathrm{H}$ (1+2+3), on $f : \cE \longrightarrow \cC$, for $\cE$ to be accessible and $f$ accessible, assuming $\cC$ \textit{Karoubian}. Or again, serviceable necessary and
sufficient conditions (stronger than 4.15.4) on $f : \cE \longrightarrow \cC$, so that every \textit{small} % [sic] suffisants -> suffisantes, agreement with "conditions"
\textit{Karoubian} category $\cC'$, and every functor $\cC' \longrightarrow \cC$, the category $\cE' = \cE \times_{\cC} \cC'$ be accessible, and the functor $\cE' \longrightarrow \cC'$ accessible. A natural idea is that it would suffice for this conclusion to be true when $\cC'$ is one of the two categories $\mathbf{\Delta}_1$, $P$ (where $P$ is the category with two objects, with three nonidentity arrows, defined above). But I doubt that this is so (\footnote[178]{That should work on the other hand if we take the categories $0 \longrightarrow 1$,
\[
\begin{array}{c}
2 \longrightarrow \underbrace{\begin{tikzcd}[ampersand replacement=\&]
1 \arrow[r, shift right, "\sigma"'] \& 0 \arrow[l, shift right, "\rho"'] \arrow[loop right, "p"]
\end{tikzcd}}_{P}
\quad \text{and} \quad
2 \longleftarrow \underbrace{\begin{tikzcd}[ampersand replacement=\&]
1 \arrow[r, shift right, "\sigma"'] \& 0 \arrow[l, shift right, "\rho"'] \arrow[loop right, "p"]
\end{tikzcd}}_{P}
\qquad (\text{\footnotemark[179]}) \ , \\
\nwarrow \ \begin{array}{c} \text{3 objects and} \\ \text{5 nonidentity} \\ \text{arrows} \end{array} \ \nearrow
\end{array}
\]
thanks to 4.15.7, page 440.}\footnotetext[179]{these categories represent respectively the functors $\mathrm{Cat} \longrightarrow \Ens$ assigning to a category $\cC$ the diagrams in $\cC$ consisting of $(x, y, p, u)$, with $p$ a projector in $p$ \texttt{[rather in $\cC$]}, and $u : y \longrightarrow x^{\natural}$ (i.e. $u : y \longrightarrow x$ with $u = pu$), resp. $u : x^{\natural} \longrightarrow y$ (i.e. $u : x \longrightarrow y$ with $u = up$).}).

We will all the same try a somewhat very special statement in this direction:

% Verified on PDF leaf8: a non-textual wavy marginal mark beside items a), b); no mathematical token omitted.
\msp{page 435}\textbf{Proposition 4.15.6} (\footnote[180]{\textit{False?}}). \textit{Let $f : \cE \longrightarrow \cC$ be a functor of $\mathfrak{U}$-categories. We assume:}
\begin{enumerate}
\item[\textit{a)}] \textit{$f$ is fibrant or cofibrant} (\footnote[181]{The fibrant case is not proved, perhaps \textit{false} -- and likewise for the cofibrant case!}).
\item[\textit{b)}] \textit{The fiber categories $\cE_s$ ($s \in \Ob \cC$) are accessible, and for every arrow $\lambda : s \longrightarrow t$ in $\cC$, the base change functor $\lambda^* : \cE_t \longrightarrow \cE_s$ (resp. cobase change $\lambda_* : \cE_s \longrightarrow \cE_t$) is accessible.}
\item[\textit{c)}] \textit{$\cC$ is essentially small and Karoubian.}
\end{enumerate}
\textit{Then $\cE$ is accessible, and $f$ is accessible.}

\textsc{Proof}. We have already seen that conditions a) and b) imply $\mathrm{H}'$, i.e. H (1+2+3). Besides, if $\cC'$ is a category and $\cC' \xrightarrow{\ \sim\ } \cC$ an equivalence, so then $\cE' \overset{\text{def}}{=} \cE \times_{\cC} \cC' \longrightarrow \cE$ is an equivalence, so the desired conclusion for $\cE$, $f$ is equivalent to the same conclusion for $\cE'$, $f' : \cE' \longrightarrow \cC'$. So we can assume that $\cC$ is small (and even reduced, if one insists). By virtue of 4.15.2 a), we thus see that the inclusion functors $\cE_s \longrightarrow \cE$ are accessible (and even $\pi_1$-accessible) as soon as $\pi_1 \geqslant \sup_{\lambda \in \Fl \cC} \pi_{\lambda}$. We also see

\msp{page 436}by AQT \texttt{[âne qui trotte: at a donkey's trot, i.e. by routine verification]} that we can find cardinals $\pi_0$, $\pi_1$ with the following properties:

\printedpage{167}
% head: A. Grothendieck : Dérivateurs, Ch. XVIII, Édition des Universités de Montpellier II et Paris VII
\begin{enumerate}
\item[\circled{1\textsuperscript{o}}] $\pi_0$ is useful for $\cE_s$, $\forall \ s \in \Ob \cC$.
\item[$2^{\circ}$)] $\pi_1 \geqslant \sup(\pi_0, \card \cC, \sup_{\lambda \in \Fl \cC} \pi_{\lambda}, \sup_{s \in \Ob \cC} \card \red \cE_s)$.
\item[\circled{3\textsuperscript{o}}] $\forall \ (\lambda : s \longrightarrow t) \in \Fl \cC$ and $x_0 \in \Ob \cE_{s\,\pi_0}$, the functor
\[
H_{\lambda}(x_0) \ : \ (x_1 \longmapsto \Hom_{\lambda}(x_0, x_1)) \ : \ \cE_t \ \longrightarrow \ (\Ens)
\]
is $\pi_1$-accessible. (\textbf{N.B.} If $f$ is cofibrant, this means that $\forall \ \lambda \in \Fl \cC$, $\lambda : s \longrightarrow t$,
\fbox{$\lambda_*(\cE_{s\,\pi_0}) \ \subset \ \cE_{s\,\pi_1}$}. If $f$ is fibrant, it suffices to require that $\forall \ \lambda \in \Fl \cC$, $\lambda^* : \cE_t \longrightarrow \cE_s$ % [sic] $\cE_{s\,\pi_1}$ -> $\cE_{t\,\pi_1}$, $\lambda_*$ goes to $\cE_t$, cf. the later display "(contenu dans 3°)"
be $\pi_1$-accessible.) But in this case, we also want \fbox{$\lambda^*(\cE_{s\,\pi_0}) \ \subset \ \cE_{s\,\pi_1}$}. % [sic] $\lambda^*(\cE_{s\,\pi_0})$ -> $\lambda^*(\cE_{t\,\pi_0})$, $\lambda^*$ is defined on $\cE_t$, cf. the later display "(contenu dans 3°)"
\item[\circled{4\textsuperscript{o}}] The functors $\lambda^*$ (case $f$ fibrant) resp. $\lambda_*$ (case $f$ cofibrant) are $\pi_1$-accessible.
\end{enumerate}
Let $\pi$ then be a cardinal such that
\[
\pi \ = \ \pi^{\pi_0} \ \geqslant \ \pi_1 \ , \tag*{$(*)$}
\]
and let $\cE(\pi)$ be the (strictly) full subcategory of $\cE$ defined by
\[
\Ob \cE(\pi) \ = \ \bigcup_{s \, \in \, \Ob \cC} \Ob \cE_{s\,\pi} \ .
\]
Thanks to $(*)$, and to $1^{\circ}$, $4^{\circ}$ and to
\[
\left\{
\begin{array}{l}
\pi_1 \ \geqslant \ \sup_s(\card \red \cE_{s\,\pi_0}) \qquad \text{(contained in $2^{\circ}$)} \\[1ex]
\lambda^*(\cE_{t\,\pi_0}) \ \subset \ \cE_{s\,\pi_1} \qquad \text{resp.} \qquad \lambda_*(\cE_{s\,\pi_0}) \ \subset \ \cE_{t\,\pi_1} \qquad \text{(contained in $3^{\circ}$)} \ ,
\end{array}
\right.
\]
we see that $\cE(\pi)$ is a

\msp{page 437}fibered subcategory, resp. a cofibered subcategory, over $\cC$ (the functor $\cE(\pi) \longrightarrow \cE$ thus being cartesian resp. cocartesian). Up to $\cC$-equivalence, $\cE$ is then reconstructed from $\cE(\pi)$, by passing to the $\Ind_{\pi}(\cE_{s\,\pi})$, which form a fibered resp. cofibered system over $\cC$ via the functors $\Ind_{\pi}(\lambda^*_{\pi})$ resp. $\Ind_{\pi}(\lambda^{\pi}_*)$ \dots

Let us prove that $\cE$ satisfies $\mathrm{L}_{\pi}$ and that $f : \cE \longrightarrow \cC$ is $\pi$-accessible (\footnote[182]{Here we use only $\pi \geqslant \card \red \cC$.}). Since $\cC$ is Karoubian, and $\pi \geqslant \pi_1 \geqslant \card \cC$, $\cC$ satisfies $\mathrm{L}_{\pi}$ and the ind-objects of $\cC$ large relative to $\pi$ are essentially constant (4.8.13). Let $I$ then be large relative to $\pi$, and $I \longrightarrow \cE$, $i \longmapsto x_i$, be a functor, we have to construct a $\varinjlim$ in $\cE$, using the already known $\varinjlim$ $s$ of the $s_i = f(x_i)$. We will take an $x$ in $\cE_s$. Let us first take the case where $f$ is cofibrant. We see that the $\beta_i(x_i)$ form a direct system in $\cE_s$, since for $i = j$, we have
\[
\beta_{i\,*}(x_j) \ = \ \beta_{j\,*} \underbrace{\beta_{ji\,*}(x_i)}_{\longrightarrow \, x_j} \ \longrightarrow \ \beta_{j\,*}(x_j) \ . % [sic] $\beta_{i\,*}(x_j)$ -> $\beta_{i\,*}(x_i)$, the object transported from $s_i$ is $x_i$
\]
\[
\begin{tikzcd}[row sep=small]
x_i \arrow[rr, "\alpha_i"] & & x \, ? \\
s_i \arrow[rr, "\beta_i"] \arrow[dr, "\beta_{ji}"'] & & s \\
& s_j \arrow[ur, "\beta_j"'] &
\end{tikzcd}
\]
We set

\printedpage{168}
% head: A. Grothendieck : Dérivateurs, Ch. XVIII, Édition des Universités de Montpellier II et Paris VII
\msp{page 438}
\[
x \ = \ \varinjlim_{I} \beta_{i\,*}(x_i) \ , \qquad \varinjlim \ \textit{in } \cE_s \ .
\]
We know that this is a $\varinjlim$ in $\cE$, hence that for every $y$ in $\cE$, we have
\[
\Hom(x, y) \xrightarrow{\ \sim\ } \varprojlim_{I} \Hom(\beta_{i\,*}(x_i), y) \ ,
\]
or what amounts to the same thing, for every arrow $\lambda : s \longrightarrow t$, and $y$ in $\Ob \cE_t$, we have
\[
\underbrace{\Hom_{\lambda}(x, y)}_{= \, \Hom_{\cE_t}(\lambda_* x, y)} \ \simeq \ \varprojlim \Hom_{\cE_t}(\underbrace{\lambda_* \beta_{i\,*}(x_i)}_{(\lambda_i \beta_i)_*(x_i)}, y) \ .
\]
This also follows from the fact that $\lambda_*$ is $\pi$-accessible (being $\pi_1$-accessible) \dots{} We easily check that this is a $\varinjlim$ of the $x_i$ in $\cE$. And of course, this limit commutes with $f$ by construction.

But in the case where $f$ is fibrant, I do not see how to go about constructing a limit in $\cE_s$. It is true that we can assume that all the $s_i$ are equal to $s$, but the transitions $\beta_{ji}$ are not the identity, nor even isomorphisms. We would have to take $\cC = P$ to try to understand -- I am not going to insist (\footnote[183]{I have checked that if $f$ is fibrant, then $\cE$ is \textit{Karoubian}. So there is hope \dots}).

\msp{page 439}In the case where $f$ is fibrant, let us admit (as an additional hypothesis if need be) that $\cE$ satisfies $L$ and that $f$ is accessible, and take $\pi$ large enough for $\cE$, $f$ to satisfy $\mathrm{L}_{\pi}$. We can then finish the proof of 4.15.6. We have to prove
\[
\cE(\pi) \ \subset \ \cE_{\pi} \ ,
\]
hence that if $x \in \Ob \cE_{s\,\pi}$, and $(y_i)_{i \in I}$ is a direct system in $\cE$ large relative to $\pi$, with limit $y$, we have
\[
\varinjlim \Hom_{\cE}(x, y_i) \xrightarrow{\ \sim\ } \Hom_{\cE}(x, y) \ .
\]
So we have to prove that every arrow $x \xrightarrow{u} y$ factors, in an essentially unique way $u_i$, through a $y_i$. But the arrow $s \xrightarrow{\ \lambda \, = \, f(u)\ } t$ in $\cC$ factors, in an essentially unique way, through some $\lambda_i \, : \, s \longrightarrow t_i$. So we are reduced to finding a \textit{$\lambda_i$-morphism} $x \longrightarrow y_i$ for large $i$, which factors. In the fibrant case, this means that we are looking for $x \longrightarrow \lambda^*_i(y_i)$ which factors $x \longrightarrow \lambda^*(y)$, now the $\lambda^*_i(y_i)$ form a direct system in $\cE_s$, and its limit \textit{in} $\cE_s$ is $\lambda^*(y)$ (\textit{to be checked}). Since $x$ is in $\cE_{s\,\pi}$, the given hom\texttt{[omorphism]} $x \longrightarrow \lambda^*(y)$ factors in

\msp{page 440}an essentially unique way through a $\lambda^*_i(y_i)$.

This time, however, it is the case where $f$ is cofibrant that poses a problem. In fact, I cannot work it out.

Thus, we do find a functor
\[
\Ind_{\pi}(\cE(\pi)) \ \longrightarrow \ \cE \ ,
\]

\end{document}
